\documentclass[aspectratio=169]{beamer}
\usetheme[noprogress]{SwissAstro}

\usepackage[french]{babel}
\usepackage[T1]{fontenc}
\usepackage[utf8]{inputenc}
\usepackage{booktabs}
\usepackage{tabularx}
\usepackage{array}
\usepackage{tikz}
\usepackage{hyperref}
\usepackage{xcolor}
\usetikzlibrary{arrows.meta,positioning,calc,shapes.geometric}

\hypersetup{colorlinks=true,urlcolor=SwissAccent,linkcolor=SwissInk}
\graphicspath{{assets/}}

\title[Programmes, défis et récompenses]{Programmes, défis et récompenses\\en astronomie amateur}
\subtitle{Du Québec à la scène internationale — choisir une quête et préparer ses premières étapes}
\author[Steve Prud'Homme]{Steve Prud'Homme}
\institute{Conférence d'astronomie amateur}
\date{Version 1.1 / SwissAstro v0.4.0 / septembre 2026}
\SwissAstroKicker{ASTRONOMIE AMATEUR / OBSERVATION / PROGRESSION}
\SwissAstroSeries{STEVE PRUD'HOMME}
\SwissAstroTitleImage{assets/demo-space.png}

% ---------------------------------------------------------------------------
% Petits composants propres à cette conférence
% ---------------------------------------------------------------------------
\newcommand{\MicroSource}[1]{\vspace{1mm}{\tiny\color{SwissMuted}#1\par}}
\newcommand{\MiniLabel}[1]{\textcolor{SwissAccent}{\bfseries\scriptsize\MakeUppercase{#1}}}

\newcommand{\FamilyCard}[3]{%
  \begin{beamercolorbox}[wd=\linewidth,sep=2.6mm,leftskip=.4mm,rightskip=.4mm]{swiss card}
    {\color{SwissAccent}\bfseries\large #1\par}
    \vspace{.8mm}
    {\bfseries #2\par}
    \vspace{.8mm}
    {\small\color{SwissMuted}#3\par}
  \end{beamercolorbox}
}

\newcommand{\ProfileCard}[4]{%
  \begin{beamercolorbox}[wd=\linewidth,sep=2.6mm,leftskip=.4mm,rightskip=.4mm]{swiss card}
    {\color{SwissAccent}\bfseries #1\par}
    {\small #2\par}
    \vspace{1mm}
    {\scriptsize\color{SwissMuted}#3\par}
    \vspace{1mm}
    {\small\bfseries #4\par}
  \end{beamercolorbox}
}

\begin{document}

% 00:00-00:03
\begin{frame}[plain]
  \titlepage
\end{frame}

% 00:03-00:07
\begin{frame}{À la fin, vous aurez choisi une prochaine quête}
\framesubtitle{Le niveau visé va de l'identification à une première mise en application.}
  \SwissAstroProcessThree
    {Identifier}{Reconnaître les principales familles de programmes, défis et récompenses.}
    {Associer}{Relier un programme à un intérêt, un équipement et un niveau d'expérience.}
    {Préparer}{Repartir avec un programme réaliste et les premières actions pour le commencer.}
  \vspace{1mm}
  \SwissAstroTakeaway{Taxonomie pédagogique : \textbf{Comprendre} $\rightarrow$ \textbf{Appliquer}. Inutile de maîtriser tous les catalogues aujourd'hui.}
\end{frame}

% 00:07-00:11
\begin{frame}{Votre pratique actuelle change le bon point de départ}
\framesubtitle{Diagnostic rapide — levez la main pour la ou les réponses qui vous ressemblent.}
  \begin{columns}[T,onlytextwidth]
    \column{0.18\textwidth}\SwissAstroBigNumber{A}{œil nu}
    \column{0.18\textwidth}\SwissAstroBigNumber{B}{jumelles}
    \column{0.18\textwidth}\SwissAstroBigNumber{C}{télescope}
    \column{0.18\textwidth}\SwissAstroBigNumber{D}{imagerie}
    \column{0.18\textwidth}\SwissAstroBigNumber{E}{science citoyenne}
  \end{columns}
  \vspace{5mm}
  \SwissAstroLead{Deuxième question : qu'est-ce qui vous motive le plus?}
  \begin{columns}[T,onlytextwidth]
    \column{0.22\textwidth}\small Apprendre le ciel
    \column{0.22\textwidth}\small Compléter une liste
    \column{0.22\textwidth}\small Contribuer à la science
    \column{0.22\textwidth}\small Créer / être reconnu
  \end{columns}
\end{frame}

\section{Donner une direction à ses nuits}

% 00:11-00:15
\begin{frame}{Un programme transforme « regarder » en « observer avec intention »}
  \begin{columns}[T,onlytextwidth]
    \column{0.43\textwidth}
      \SwissAstroTag{SANS CADRE}
      \vspace{2mm}
      {\fontsize{18}{20}\selectfont\bfseries « Qu'est-ce que je regarde ce soir? »}\par
      \vspace{3mm}
      {\small\color{SwissMuted}On revient facilement aux mêmes cibles, sans trace claire de sa progression.}
    \column{0.49\textwidth}
      \SwissAstroTag{AVEC UN PROGRAMME}
      \vspace{2mm}
      {\fontsize{18}{20}\selectfont\bfseries Objectif + méthode + journal + validation}\par
      \vspace{3mm}
      {\small\color{SwissMuted}La liste devient un parcours d'apprentissage plutôt qu'une simple collection de coches.}
  \end{columns}
  \vspace{2mm}
  \SwissAstroTakeaway{La récompense est visible; la compétence acquise est la vraie valeur du parcours.}
\end{frame}

% 00:15-00:19
\begin{frame}{Il existe quatre grandes familles de défis}
\framesubtitle{Cette carte servira de repère pendant toute la conférence.}
  \SwissAstroFourColumns{
    \FamilyCard{01}{Certification par liste}{Observer une liste prescrite et documenter ses observations.}
  }{
    \FamilyCard{02}{Contribution scientifique}{Produire et transmettre des données utilisables par la recherche.}
  }{
    \FamilyCard{03}{Concours / mérite}{Faire reconnaître une image, un projet, un service ou une réalisation.}
  }{
    \FamilyCard{04}{Défi communautaire}{Participer à un défi ponctuel, mensuel ou événementiel.}
  }
  \MicroSource{Synthèse du document fourni; exemples : SRAC, Astronomical League, AAVSO, IOTA, FAAQ et concours d'astrophotographie.}
\end{frame}

% 00:19-00:25
\begin{frame}{Activité 1 — Classer avant de mémoriser}
\framesubtitle{Classez A à F dans les quatre familles — 3 minutes, puis correction.}
  \begin{columns}[T,onlytextwidth]
    \column{0.47\textwidth}
      \MiniLabel{A — 110 objets Messier}\par
      {\small Observer, consigner, faire valider une liste.}\par\vspace{2.6mm}
      \MiniLabel{B — étoiles variables}\par
      {\small Mesurer des magnitudes et soumettre les observations.}\par\vspace{2.6mm}
      \MiniLabel{C — concours photo}\par
      {\small Soumettre une image selon des critères annoncés.}
    \column{0.45\textwidth}
      \MiniLabel{D — défi mensuel EAA}\par
      {\small Nouvelles cibles proposées par une communauté.}\par\vspace{2.6mm}
      \MiniLabel{E — Trophée Méritas}\par
      {\small Reconnaissance d'un engagement soutenu.}\par\vspace{2.6mm}
      \MiniLabel{F — occultation d'astéroïde}\par
      {\small Chronométrer un événement et transmettre la donnée.}
  \end{columns}
\end{frame}

% 00:25-00:27
\begin{frame}{Correction — ce qui compte, c'est la logique du programme}
  \begin{columns}[T,onlytextwidth]
    \column{0.22\textwidth}\SwissAstroBigNumber{A}{liste}\vspace{2mm}\SwissAstroBigNumber{B}{science}
    \column{0.22\textwidth}\SwissAstroBigNumber{C}{concours}\vspace{2mm}\SwissAstroBigNumber{D}{communauté}
    \column{0.22\textwidth}\SwissAstroBigNumber{E}{mérite}\vspace{2mm}\SwissAstroBigNumber{F}{science}
    \column{0.22\textwidth}
      \SwissAstroTag{REPÈRE}
      \vspace{2mm}
      {\small Une même pratique peut croiser plusieurs familles. Ce sont les \textbf{règles de réussite} qui déterminent la famille.}
  \end{columns}
\end{frame}

\section{Québec — une progression proche de chez nous}

% 00:27-00:31
\begin{frame}{La FAAQ propose maintenant 11 programmes d'observation}
\framesubtitle{Mise à jour officielle vérifiée le 10 septembre 2026.}
  \begin{columns}[T,onlytextwidth]
    \column{0.30\textwidth}
      \SwissAstroBigNumber{11}{programmes}
      \vspace{3mm}
      {\small\color{SwissMuted}Plusieurs portes d'entrée ne demandent pas d'équipement spécialisé.}
    \column{0.62\textwidth}
      \begin{columns}[T,onlytextwidth]
        \column{0.48\textwidth}\small
          \textbf{Lune}\\ Soleil\\ Rayonnement\\ Système solaire\\ Astrophotographie\\ Herschel
        \column{0.48\textwidth}\small
          \textbf{Étoiles doubles}\\ Constellations\\ Caldwell\\ Messier\\ 50e anniversaire
      \end{columns}
      \vspace{2mm}
      \SwissAstroTakeaway{Le paysage québécois est plus structuré aujourd'hui que ne le laissait entendre le document de départ.}
  \end{columns}
\end{frame}

% 00:31-00:35
\begin{frame}{La progression FAAQ se lit en quatre niveaux}
  \begin{columns}[T,onlytextwidth]
    \column{0.21\textwidth}\SwissAstroBigNumber{1}{Débutant}\small accessible à tous, sauf Rayonnement
    \column{0.21\textwidth}\SwissAstroBigNumber{2}{Curieux}\small abonnement FAAQ / club affilié
    \column{0.21\textwidth}\SwissAstroBigNumber{3}{Motivé}\small exigences plus soutenues
    \column{0.21\textwidth}\SwissAstroBigNumber{4}{Expert}\small maîtrise et persévérance
  \end{columns}
  \vspace{5mm}
  \SwissAstroTakeaway{Chaque niveau réussi donne droit à une \textbf{épinglette}. Les programmes Constellations et 50e anniversaire font exception à la structure à quatre niveaux.}
  \MicroSource{FAAQ — Les programmes d'observation, 2026.}
\end{frame}

% 00:35-00:38
\begin{frame}{Trois portes d'entrée donnent rapidement une première victoire}
\framesubtitle{Le bon début n'est pas forcément le programme le plus prestigieux.}
  \SwissAstroOpenThree
    {Constellations}{Développer le repérage du ciel sans dépendre d'un gros instrument.}
    {Lune}{Observer souvent, depuis la ville, et apprendre à décrire des détails.}
    {Messier}{Passer progressivement au ciel profond avec une liste classique et motivante.}
  \vspace{3mm}
  \SwissAstroTakeaway{Choisissez d'abord un programme compatible avec votre ciel, votre temps et votre équipement réel.}
\end{frame}

% 00:38-00:41
\begin{frame}{La FAAQ reconnaît aussi les pratiques modernes}
\framesubtitle{Le cadre d'apprentissage peut intégrer les télescopes intelligents.}
  \begin{columns}[T,onlytextwidth]
    \column{0.44\textwidth}
      \SwissAstroTag{NIVEAU DÉBUTANT}
      \vspace{2mm}
      {\fontsize{19}{21}\selectfont\bfseries Seestar, Vespera, eVscope...}\par
      \vspace{2mm}
      {\small Certains programmes autorisent ces instruments au niveau débutant.}
    \column{0.48\textwidth}
      \SwissAstroTag{ASTROPHOTOGRAPHIE}
      \vspace{2mm}
      {\fontsize{19}{21}\selectfont\bfseries Planifier + acquérir + traiter + présenter}\par
      \vspace{2mm}
      {\small La compétence évaluée est un processus, pas seulement une belle image finale.}
  \end{columns}
  \MicroSource{FAAQ — Les programmes d'observation et programme Astrophotographie, consultés en septembre 2026.}
\end{frame}

% 00:41-00:45
\begin{frame}{Au Québec, la reconnaissance dépasse la complétion de listes}
  \SwissAstroOpenFour
    {Méritas}{Engagement soutenu envers la Fédération ou un club.}
    {Les Pléiades}{Contribution remarquable d'une personne de moins de 20 ans.}
    {Qilak}{Vulgarisation bénévole exceptionnelle auprès des jeunes et du public.}
    {20+ programmes}{Reconnaissance spéciale pour une participation exceptionnelle aux programmes d'observation.}
  \MicroSource{FAAQ — Reconnaissance au mérite, consulté le 10 septembre 2026.}
\end{frame}

% 00:45-00:48
\begin{frame}{CAFTA — fabriquer, rechercher, imager}
\framesubtitle{Exemple québécois où l'excellence technique devient une voie de reconnaissance.}
  \begin{columns}[T,onlytextwidth]
    \column{0.33\textwidth}
      \SwissAstroBigNumber{43e}{CAFTA 2026}
      \vspace{2mm}
      {\small 17 octobre 2026 / Dorval}
    \column{0.59\textwidth}
      \SwissAstroLead{Catégories annoncées}
      \begin{itemize}
        \item Finesse du travail
        \item Recherche et technologie
        \item Imagerie avancée
        \item Prix Fred-Clarke
      \end{itemize}
  \end{columns}
  \MicroSource{FAAQ — CAFTA 2026, page officielle de l'événement.}
\end{frame}

\begin{frame}{Vérification express — Québec}
\framesubtitle{Répondez sans vos notes.}
  \SwissAstroThreeColumns{
    \SwissAstroNumberedCard{01}{Combien?}{Combien de programmes d'observation la FAAQ affiche-t-elle actuellement?}
  }{
    \SwissAstroNumberedCard{02}{Progression}{Nommez deux des quatre niveaux d'expertise.}
  }{
    \SwissAstroNumberedCard{03}{Reconnaissance}{Citez une récompense qui ne dépend pas simplement d'une liste d'objets.}
  }
  \vspace{4mm}
  {\small\color{SwissMuted}Réponses : 11; Débutant / Curieux / Motivé / Expert; Méritas, Pléiades, Qilak, CAFTA, etc.}
\end{frame}

\section{Canada — la progression structurée de la SRAC}

% 00:48-00:52
\begin{frame}{La SRAC offre huit grands programmes visuels certifiants}
  \begin{columns}[T,onlytextwidth]
    \column{0.31\textwidth}
      \SwissAstroBigNumber{08}{programmes visuels}
      \vspace{3mm}
      {\small Plusieurs mènent à un certificat et à une épinglette.}
    \column{0.61\textwidth}\small
      \begin{columns}[T,onlytextwidth]
        \column{0.48\textwidth}
          Explore the Universe\\ Explore the Moon\\ Messier Catalogue\\ Finest NGC Objects
        \column{0.48\textwidth}
          Isabel Williamson Lunar\\ Deep-Sky Gems\\ Deep-Sky Challenge\\ Double Stars
      \end{columns}
      \vspace{4mm}
      {\small\color{SwissMuted}À côté de ces huit parcours visuels, la SRAC présente aussi un certificat d'astro-imagerie et des projets d'observation spéciaux.}
  \end{columns}
  \MicroSource{RASC/SRAC — rasc.ca/certificate-programs et rasc.ca/observing, consultés en 2026.}
\end{frame}

% 00:52-00:56
\begin{frame}{Pour débuter, la SRAC réduit volontairement la barrière d'entrée}
  \SwissAstroTwoColumns[0.44][0.48]{
    \SwissAstroTag{EXPLORE THE UNIVERSE}
    \vspace{2mm}
    \SwissAstroLead{Novice / œil nu / jumelles / petit instrument}
    {\small Couvre constellations, étoiles, Lune, Système solaire, ciel profond et étoiles doubles. Ouvert aussi aux non-membres.}
  }{
    \SwissAstroTag{EXPLORE THE MOON}
    \vspace{2mm}
    \SwissAstroLead{100 formations lunaires}
    {\small Deux voies de reconnaissance selon l'instrument : jumelles ou télescope. Un terrain idéal pour apprendre à décrire précisément ce qu'on voit.}
  }
  \MicroSource{RASC — Visual Observing Programs et Explore the Universe.}
\end{frame}

% 00:56-01:00
\begin{frame}{Le niveau intermédiaire permet de choisir une spécialité}
  \SwissAstroOpenFour
    {Messier}{110 objets; classique du ciel profond.}
    {Finest NGC}{110 objets; étape plus exigeante en ciel profond.}
    {I. Williamson}{Parcours lunaire beaucoup plus détaillé.}
    {Double Stars}{110 systèmes doubles ou multiples.}
  \vspace{3mm}
  \SwissAstroTakeaway{Le choix intermédiaire doit suivre ce que vous aimez réellement observer — pas une hiérarchie sociale des instruments.}
  \MicroSource{RASC — Visual Observing Programs, 2026.}
\end{frame}

% 01:00-01:03
\begin{frame}{Deux parcours poussent le ciel profond jusqu'au niveau avancé}
  \SwissAstroCompare{DEEP-SKY GEMS}{
    \SwissAstroBigNumber{154}{objets}
    \vspace{2mm}
    {\small Sélection de joyaux du ciel profond; certificat.}
  }{DEEP-SKY CHALLENGE}{
    \SwissAstroBigNumber{45}{objets difficiles}
    \vspace{2mm}
    {\small Programme conçu pour des observations nettement plus exigeantes; certificat.}
  }
  \MicroSource{RASC — Visual Observing Programs; document de synthèse fourni.}
\end{frame}

% 01:03-01:07
\begin{frame}{Le journal d'observation est la compétence transversale}
\framesubtitle{Tous les programmes certifiants visuels de la SRAC exigent un journal.}
  \begin{columns}[T,onlytextwidth]
    \column{0.32\textwidth}
      \SwissAstroBigNumber{01}{objet}
      \vspace{2mm}\SwissAstroBigNumber{02}{conditions}
      \vspace{2mm}\SwissAstroBigNumber{03}{description}
    \column{0.60\textwidth}
      \SwissAstroCard{Entrée minimale utile}{Date et heure / lieu / instrument / grossissement / conditions / objet / ce qui est réellement vu.}
      \vspace{2mm}
      \SwissAstroCard{Règle pédagogique}{Localiser soi-même, observer soi-même, tenir son propre journal et soumettre son dossier sur son propre mérite.}
  \end{columns}
  \MicroSource{RASC — Visual Observing Programs, sections Log Requirements et Individual Effort.}
\end{frame}

% 01:07-01:10
\begin{frame}{GoTo ou repérage manuel : le programme doit préciser la règle}
  \begin{columns}[T,onlytextwidth]
    \column{0.44\textwidth}
      \SwissAstroTag{DEUX VERSIONS}
      \vspace{2mm}
      {\small\bfseries Messier / Finest NGC / Double Stars}
      \vspace{2mm}
      {\small La SRAC distingue des voies traditionnelles et assistées par ordinateur pour certains certificats.}
    \column{0.48\textwidth}
      \SwissAstroTag{À RETENIR}
      \vspace{2mm}
      {\fontsize{17}{19}\selectfont\bfseries Lire les règles \emph{avant} la première observation.}
      \vspace{2mm}
      {\small Un équipement permis dans un programme peut être interdit ou traité différemment dans un autre.}
  \end{columns}
  \MicroSource{RASC — Visual Observing Programs; détails des certificats.}
\end{frame}

% 01:10-01:18
\begin{frame}{Activité 2 — Quel programme pour quel profil?}
\framesubtitle{4 minutes en petits groupes + 4 minutes de retour. Un critère décisif par choix.}
  \SwissAstroThreeColumns{
    \SwissAstroCard{Alex}{\textbf{Jumelles 10×50 / débutant / ville}\par Veut apprendre le ciel sans acheter d'autre matériel.\par\vspace{1mm}\textcolor{SwissAccent}{\bfseries Quel programme SRAC?}}
  }{
    \SwissAstroCard{Myriam}{\textbf{Dobson 200 mm / Messier déjà amorcé}\par Veut progresser en ciel profond de façon structurée.\par\vspace{1mm}\textcolor{SwissAccent}{\bfseries Quel programme SRAC?}}
  }{
    \SwissAstroCard{Karim}{\textbf{Télescope 90 mm / aime les objets serrés}\par Veut développer sa finesse d'observation.\par\vspace{1mm}\textcolor{SwissAccent}{\bfseries Quel programme SRAC?}}
  }
  \vspace{2mm}
  {\small\bfseries Cherchez l'adéquation objectif–instrument–difficulté, pas un « meilleur » programme absolu.}
\end{frame}

% 01:18-01:20
\begin{frame}{Correction — trois bonnes portes d'entrée, pour trois raisons différentes}
  \SwissAstroThreeColumns{
    \SwissAstroCard{Alex → Explore the Universe}{Accessible au novice et aux non-membres; œil nu, jumelles ou petit instrument.}
  }{
    \SwissAstroCard{Myriam → Finest NGC}{Progression logique après Messier pour un observateur de ciel profond bien équipé.}
  }{
    \SwissAstroCard{Karim → Double Stars}{Programme spécialisé compatible avec un petit télescope; développe l'attention aux détails.}
  }
  \vspace{4mm}
  \SwissAstroTakeaway{Associer correctement un programme à un profil est déjà une compétence d'\textbf{application}.}
\end{frame}

\begin{SwissAstroDarkFrame}
  \vfill
  \SwissAstroTag{PAUSE / 10 MINUTES}
  \vspace{7mm}
  {\fontsize{30}{33}\selectfont\bfseries\color{SwissWhite}Revenez avec le nom d'un programme que vous pourriez réellement commencer.}\par
  \vspace{6mm}
  {\small\color{SwissRule}À la reprise : États-Unis, science citoyenne, occultations et astrophotographie internationale.}
  \vfill
\end{SwissAstroDarkFrame}

\section{International — spécialiser, contribuer, créer}

% 01:30-01:34
\begin{frame}{L'international élargit surtout trois choses : choix, science et visibilité}
  \SwissAstroProcessThree
    {Se spécialiser}{Des dizaines de programmes très ciblés via l'Astronomical League.}
    {Contribuer}{Étoiles variables, occultations et autres données utiles à la recherche.}
    {Créer}{Concours d'astrophotographie à portée francophone ou mondiale.}
\end{frame}

% 01:34-01:38
\begin{frame}{Astronomical League : plus de 75 programmes d'observation}
\framesubtitle{Une organisation américaine, mais un catalogue de programmes à portée internationale.}
  \begin{columns}[T,onlytextwidth]
    \column{0.32\textwidth}
      \SwissAstroBigNumber{75+}{programmes}
      \vspace{3mm}
      {\small\color{SwissMuted}La plupart des reconnaissances demandent d'être membre par une société affiliée ou comme Member-at-Large.}
    \column{0.60\textwidth}
      \SwissAstroLead{Exemples de niches}
      \begin{columns}[T,onlytextwidth]
        \column{0.48\textwidth}\small Constellations\\ Lune\\ Messier\\ Jumelles\\ Herschel
        \column{0.48\textwidth}\small Astéroïdes\\ Étoiles carbonées\\ Galaxies\\ Spectroscopie\\ Imagerie
      \end{columns}
  \end{columns}
  \MicroSource{Astronomical League — Master Observer Progression et pages de programmes, consultées en 2026.}
\end{frame}

% 01:38-01:41
\begin{frame}{Master Observer — une progression sur plusieurs années}
\framesubtitle{La récompense relie plusieurs programmes en un parcours à long terme.}
  \centering
  \begin{tikzpicture}[scale=.82,transform shape]
    \tikzset{box/.style={draw=SwissInk,fill=SwissPaper,rounded corners=1pt,minimum width=2.25cm,minimum height=.78cm,align=center,font=\small\bfseries},
             arrow/.style={-{Latex[length=2mm]},line width=.7pt,draw=SwissInk}}
    \node[box] (o) at (0,0) {Observer};
    \node[box] (m) at (2.8,0) {Master Observer};
    \node[box] (a) at (5.6,0) {Advanced};
    \node[box] (s) at (8.4,0) {Silver};
    \node[box] (g) at (11.2,0) {Gold};
    \node[box,fill=SwissAccent,text=white,draw=SwissAccent] (p) at (11.2,-1.5) {Platinum};
    \draw[arrow] (o)--(m); \draw[arrow] (m)--(a); \draw[arrow] (a)--(s); \draw[arrow] (s)--(g); \draw[arrow] (g)--(p);
  \end{tikzpicture}
  \vspace{3mm}
  {\small\color{SwissMuted}Branches complémentaires : Binocular Master Observer / Master Imager / For Your Eyes Only / Triple-Crown.}\par
  \vspace{3mm}
  {\small\bfseries Un objectif de carrière amateur — pas un projet à terminer en une saison.}
  \MicroSource{Astronomical League — Master Observer Progression.}
\end{frame}

% 01:41-01:45
\begin{frame}{AAVSO — la réussite se mesure par la qualité de la contribution}
  \begin{columns}[T,onlytextwidth]
    \column{0.32\textwidth}
      \SwissAstroTag{SCIENCE CITOYENNE}
      \vspace{2mm}
      \SwissAstroBigNumber{1}{code observateur}
      \vspace{2mm}
      {\small Créer un compte, demander un code et commencer avec des cibles adaptées.}
    \column{0.60\textwidth}
      \SwissAstroLead{Boucle d'apprentissage}
      \begin{enumerate}
        \item choisir une étoile variable accessible;
        \item estimer ou mesurer sa magnitude;
        \item documenter l'observation;
        \item soumettre des données utilisables;
        \item améliorer sa méthode avec la communauté et le mentorat.
      \end{enumerate}
  \end{columns}
  \MicroSource{AAVSO — Honors and Awards; outils pour observateurs et inscription, consultés en 2026.}
\end{frame}

% 01:45-01:49
\begin{frame}{IOTA — transformer une occultation en donnée scientifique}
  \begin{columns}[T,onlytextwidth]
    \column{0.54\textwidth}
      \begin{tikzpicture}[scale=.82]
        \fill[SwissInk] (0,0) circle (1.5pt);
        \node[anchor=south,font=\scriptsize] at (0,.15) {étoile};
        \fill[SwissAccent] (4,0) circle (7pt);
        \node[anchor=south,font=\scriptsize] at (4,.25) {astéroïde};
        \draw[-{Latex[length=2.5mm]},line width=1pt,SwissAccent] (5.3,0)--(2.5,0);
        \draw[line width=.7pt,SwissMuted,dashed] (0,-1.2)--(0,1.2);
        \node[align=center,font=\small] at (2,-1.25) {chronométrer disparition\\et réapparition};
      \end{tikzpicture}
    \column{0.38\textwidth}
      \SwissAstroLead{Ce qui est reconnu}
      {\small Contributions significatives à la science des occultations, aux campagnes, aux outils et à l'organisation.}\par
      \vspace{2mm}
      {\scriptsize\color{SwissMuted}Depuis février 2026, les rapports d'occultations astéroïdales sont orientés vers IOTA-Hub.}
  \end{columns}
  \MicroSource{IOTA — IOTA Awards; Reporting Observations, mises à jour 2026.}
\end{frame}

% 01:49-01:53
\begin{frame}{Astrophotographie — du Québec aux concours internationaux}
  \SwissAstroTwoColumns[0.44][0.48]{
    \SwissAstroTag{FRANCOPHONIE}
    \vspace{2mm}
    \SwissAstroLead{Les Étoiles de l'Astronomie — AFA}
    {\small Plusieurs catégories sont explicitement ouvertes aux photographes résidant dans la province du Québec.}
  }{
    \SwissAstroTag{MONDE}
    \vspace{2mm}
    \SwissAstroLead{ZWO Astronomy Photographer of the Year}
    {\small Concours mondial du Royal Observatory Greenwich; jusqu'à dix images par participant. L'édition 2026 est close; résultats en septembre.}
  }
  \vspace{2mm}
  \MicroSource{Sources : AFA — Les Étoiles de l'Astronomie; Royal Museums Greenwich — APY 18, 2026.}
\end{frame}

% 01:53-01:56
\begin{frame}{Même ciel, quatre philosophies de réussite}
  \begin{columns}[T,onlytextwidth]
    \column{0.22\textwidth}\SwissAstroCard{FAAQ}{\textbf{Communauté}\par Progresser près de chez soi et être reconnu par ses pairs.}
    \column{0.22\textwidth}\SwissAstroCard{SRAC}{\textbf{Curriculum}\par Développer progressivement des compétences d'observation.}
    \column{0.22\textwidth}\SwissAstroCard{AL}{\textbf{Spécialisation}\par Explorer une grande diversité de niches et de paliers.}
    \column{0.22\textwidth}\SwissAstroCard{AAVSO / IOTA / concours}{\textbf{Contribution ou création}\par Produire une donnée utile ou une œuvre reconnue.}
  \end{columns}
\end{frame}

\section{Choisir sa prochaine quête}

% 01:56-02:00
\begin{frame}{Le meilleur programme est celui que vous pouvez réellement commencer}
\framesubtitle{Cinq filtres suffisent pour réduire rapidement le choix.}
  \SwissAstroOpenThree
    {01 — Motivation}{Qu'est-ce qui vous donne envie de sortir observer?}
    {02 — Équipement}{Que pouvez-vous utiliser régulièrement, sans achat préalable?}
    {03 — Ciel + temps}{Quel niveau de pollution lumineuse et quelle disponibilité réelle?}
  \vspace{3mm}
  \begin{columns}[T,onlytextwidth]
    \column{0.46\textwidth}\SwissAstroOpenItem[04]{Règles}{Méthode permise, journal, validation et admissibilité.}
    \column{0.46\textwidth}\SwissAstroOpenItem[05]{Soutien}{Club, mentor ou communauté capable de vous accompagner.}
  \end{columns}
  \vspace{3mm}
  \SwissAstroTakeaway{Un programme « ambitieux » qui reste dans un signet n'enseigne rien. Un programme commencé dès ce mois-ci, oui.}
\end{frame}

% 02:00-02:03
\begin{frame}{Arbre de décision — partez de votre intention}
\framesubtitle{Trois branches simples suffisent pour trouver une première famille de programmes.}
  \SwissAstroThreeColumns{
    \SwissAstroTag{OBSERVER}
    \vspace{2mm}
    {\bfseries Apprendre / progresser}\par
    \vspace{2mm}{\small FAAQ ou SRAC pour un parcours structuré.}\par
    \vspace{3mm}\textcolor{SwissAccent}{\bfseries $\downarrow$ AL}\par
    {\scriptsize\color{SwissMuted}si vous cherchez une spécialisation très précise}
  }{
    \SwissAstroTag{CONTRIBUER}
    \vspace{2mm}
    {\bfseries Produire une donnée utile}\par
    \vspace{2mm}{\small AAVSO pour les variables; IOTA pour les occultations.}\par
    \vspace{3mm}\textcolor{SwissAccent}{\bfseries $\downarrow$ science citoyenne}
  }{
    \SwissAstroTag{CRÉER}
    \vspace{2mm}
    {\bfseries Faire reconnaître une image}\par
    \vspace{2mm}{\small AFA pour une scène francophone; ZWO APY pour une scène mondiale.}\par
    \vspace{3mm}\textcolor{SwissAccent}{\bfseries $\downarrow$ concours}
  }
\end{frame}

% 02:03-02:12
\begin{frame}{Activité 3 — Mon prochain programme}
\framesubtitle{8 à 9 minutes — un plan concret à commencer dans les 30 prochains jours.}
  \begin{columns}[T,onlytextwidth]
    \column{0.46\textwidth}
      \SwissAstroNumberedCard{1}{Je choisis}{Nom du programme / organisation / niveau visé.}
      \vspace{1.5mm}
      \SwissAstroNumberedCard{2}{Je vérifie}{Admissibilité, équipement, méthode permise, journal et preuve demandée.}
    \column{0.46\textwidth}
      \SwissAstroNumberedCard{3}{Je commence}{Mes trois premières observations ou actions concrètes.}
      \vspace{1.5mm}
      \SwissAstroNumberedCard{4}{Je m'engage}{Une première date au calendrier + une personne ou un club pour m'aider.}
  \end{columns}
\end{frame}

% 02:12-02:15
\begin{frame}{Checklist avant la première observation}
  \begin{columns}[T,onlytextwidth]
    \column{0.18\textwidth}\SwissAstroStep{01}{Règles}{Lire la version actuelle.}
    \column{0.18\textwidth}\SwissAstroStep{02}{Admissibilité}{Membre? niveau? âge?}
    \column{0.18\textwidth}\SwissAstroStep{03}{Matériel}{Instrument et méthode permis.}
    \column{0.18\textwidth}\SwissAstroStep{04}{Journal}{Créer un modèle simple.}
    \column{0.18\textwidth}\SwissAstroStep{05}{Première cible}{Planifier petit et proche.}
  \end{columns}
  \vspace{6mm}
  \SwissAstroTakeaway{Avant d'accumuler des observations, vérifiez surtout \textbf{ce qui doit être documenté et comment la réussite sera validée}.}
\end{frame}

% 02:15-02:17
\begin{frame}{Un plan de 30 jours suffit pour passer de l'intention à la pratique}
  \begin{columns}[T,onlytextwidth]
    \column{0.22\textwidth}\SwissAstroTimelineItem{Semaine 1}{Préparer}{Règles, liste, journal, logiciel ou atlas.}
    \column{0.22\textwidth}\SwissAstroTimelineItem{Semaine 2}{Observer}{Réaliser 2 à 3 observations et les consigner.}
    \column{0.22\textwidth}\SwissAstroTimelineItem{Semaine 3}{Valider}{Faire relire un exemple de journal par un pair ou un club.}
    \column{0.22\textwidth}\SwissAstroTimelineItem{Semaine 4}{Répéter}{Ajuster la méthode et fixer le prochain petit jalon.}
  \end{columns}
\end{frame}

% 02:17-02:19
\begin{frame}{Sortie 3–2–1}
\framesubtitle{Dernière récupération active avant les questions.}
  \SwissAstroThreeColumns{
    \SwissAstroBigNumber{3}{noms}\vspace{2mm}{\small Trois organisations ou programmes que vous retenez.}
  }{
    \SwissAstroBigNumber{2}{critères}\vspace{2mm}{\small Deux critères qui orientent votre choix personnel.}
  }{
    \SwissAstroBigNumber{1}{action}\vspace{2mm}{\small Une action datée pour commencer votre programme.}
  }
\end{frame}

\SwissAstroClosingFrame{Le ciel est le même. Les quêtes sont multiples.}{Questions / discussion — et surtout : quelle quête allez-vous commencer?}

% ---------------------------------------------------------------------------
% Annexes — à utiliser selon les questions ou comme matériel de référence.
% ---------------------------------------------------------------------------
\appendix
\section*{Annexes}

\begin{frame}{Ressources de départ}
  \small
  \begin{tabularx}{\linewidth}{>{\bfseries}p{0.23\linewidth} X}
    FAAQ & \href{https://www.faaq.org/les-programmes-dobservation/}{faaq.org/les-programmes-dobservation} \\
    SRAC / RASC & \href{https://www.rasc.ca/certificate-programs}{rasc.ca/certificate-programs} \\
    Astronomical League & \href{https://www.astroleague.org/master-observer-progression/}{astroleague.org/master-observer-progression} \\
    AAVSO & \href{https://www.aavso.org/}{aavso.org} \\
    IOTA & \href{https://occultations.org/}{occultations.org} \\
    AFA — concours & \href{https://www.afastronomie.fr/les-etoiles-de-l-astronomie}{afastronomie.fr/les-etoiles-de-l-astronomie} \\
    ZWO APY & \href{https://www.rmg.co.uk/whats-on/astronomy-photographer-year/competition}{rmg.co.uk / Astronomy Photographer of the Year} \\
  \end{tabularx}
  \vspace{3mm}
  {\scriptsize\color{SwissMuted}Toujours lire la page officielle du programme avant de commencer : règles, admissibilité, méthodes permises et formulaires peuvent évoluer.}
\end{frame}

\begin{frame}{SRAC — repère comparatif des huit programmes visuels}
\scriptsize
\begin{tabularx}{\linewidth}{>{\bfseries}p{0.24\linewidth} p{0.14\linewidth} p{0.15\linewidth} X}
\toprule
Programme & Niveau & Volume indicatif & Point d'attention \\
\midrule
Explore the Universe & Débutant & sélection d'objets & ouvert aux non-membres; idéal pour commencer \\
Explore the Moon & Débutant & 100 formations & voie jumelles / voie télescope \\
Messier Catalogue & Intermédiaire & 110 objets & versions traditionnelles / assistées selon règles SRAC \\
Finest NGC Objects & Intermédiaire & 110 objets & ciel profond plus exigeant \\
Isabel Williamson Lunar & Intermédiaire & parcours détaillé & observation lunaire soutenue \\
Double Stars & Intermédiaire & 110 systèmes & petit instrument possible \\
Deep-Sky Gems & Avancé & 154 objets & certificat \\
Deep-Sky Challenge & Avancé & 45 défis & plusieurs cibles difficiles \\
\bottomrule
\end{tabularx}
\MicroSource{RASC/SRAC — Visual Observing Programs; synthèse fournie. Vérifier la page de chaque programme pour les exigences exactes.}
\end{frame}

\begin{frame}{FAAQ — les 11 programmes visibles en septembre 2026}
  \SwissAstroOpenThree
    {Système solaire}{Lune\\Soleil\\Système solaire\\Rayonnement}
    {Ciel profond / repérage}{Constellations\\Messier\\Caldwell\\Herschel\\Étoiles doubles}
    {Imagerie / spécial}{Astrophotographie\\50e anniversaire}
  \vspace{2mm}
  {\small\bfseries La FAAQ indique quatre niveaux pour la plupart des programmes et une épinglette à chaque niveau réussi.}
\end{frame}

\begin{frame}{International — exemples à explorer après la conférence}
  \SwissAstroOpenFour
    {AL — fondamentaux}{Constellation Hunter\\Lunar\\Messier\\Solar System}
    {AL — spécialités}{Herschel\\Astéroïdes\\Spectroscopie\\Imagerie}
    {Science citoyenne}{AAVSO — variables\\IOTA — occultations}
    {Création}{AFA — Étoiles de l'Astronomie\\ZWO APY — Greenwich}
\end{frame}

\end{document}
